\section{Related Work}

% Cao et al. \cite{cao2015spiking} pioneered the field of ANN-to-SNN conversion, laying the foundation for subsequent research in this area . Following this work, numerous studies have focused on improving neuronal representation capabilities and optimizing spike firing characteristics. For example, Li, Wang, and colleagues introduced the value -1 into the ANN-to-SNN conversion paradigm, effectively improving the information representation capacity of neurons \cite{li2022quantization, SNM, spikeziptf}. Hao et al. ~\cite{hao2023bridging} addressed the delay in spike firing by optimizing initial membrane potential settings, thereby improving the efficiency of temporal encoding. Huang et al.~\cite{ECMT} (ECMT) proposed a multi-threshold design that further enhances neuronal representation; this approach aligns well with binary-based chip architectures, enabling computational savings (e.g., multiplication/division by two corresponds to increment/decrement operations on binary exponent bits).

% In addition, DCGS introduced a differential encoding scheme that further improves spike representation by allowing different timesteps to use threshold/T as a unit for information encoding, significantly expanding the representational capacity \cite{dcgs} (DCGS). However, the extensive use of division by
% T operations in its encoding and decoding processes may incur high energy consumption on existing binary-based hardware.

% Jiang et al.\cite{adaptive} proposed the Adaptive Fission algorithm, which reduces the latency requirements of the conversion of ANN-to-SNN by dynamically increasing the number of neurons, trading spatial complexity for temporal efficiency. Nevertheless, there remains room for improvement in both conversion accuracy and latency.

% Recent works have successfully extended ANN-to-SNN conversion techniques to large language models (LLMs) \cite{spikellm, fas, las}(spikellm, fas, las). However, due to the inherent complexity of LLMs, these approaches still face challenges in achieving satisfactory conversion accuracy with low latency.

% Building upon prior research, this work leverages the binary design principles of modern hardware and exploits the characteristics of data's binary representation. Our approach reduces computational cost while further enhancing information representation, providing a new pathway for efficient ANN-to-SNN conversion.

\subsection{Training-Dependent Conversion}
 
Training-dependent methods retrain the source ANN to align activations with spiking neuron dynamics before conversion. Under rate coding, representative works optimize membrane potential initialization~\cite{bu2022optimized}, replace ReLU with quantization-aware activations~\cite{bu2023optimal}, or apply dual-phase pipelines for weight and threshold co-calibration~\cite{ding2021optimal, wang2022towards}. For Transformers, SpikeZIP-TF~\cite{you2024spikezip} fine-tunes pretrained weights with spike-compatible surrogates, and SpikeLLM~\cite{xing2024spikellm} extends spiking mechanisms to billion-parameter LLMs via saliency-aware encoding. While effective, these methods require training pipeline access, limiting applicability when the source ANN is fixed or proprietary.

\subsection{Training-Free Conversion}
 
In contrast, Training-free approaches convert pretrained ANNs directly without fine-tuning. The dominant paradigm is rate coding~\cite{rueckauer2017conversion, diehl2015fast}, which requires at least $T = N$ timesteps to distinguish $N$ activation levels. To mitigate this, prior works calibrate residual and initial membrane potentials~\cite{hao2023reducing, hao2023bridging}, extend neuron outputs to ternary values~\cite{wang2022signed}, or refine threshold calibration strategies~\cite{wang2025adaptive, bu2025inference}. Beyond rate coding, phase coding~\cite{kim2018deep, wang2022efficient} supports $2^T$ representable values but requires all timesteps to elapse before decoding, while differential coding (DCGS)~\cite{huang2025differential} eliminates temporal decay at the cost of hardware-unfriendly division operations. For Transformer conversion, ECMT~\cite{huang2024towards} employs multi-threshold neurons but retains floating-point multiplication; STA~\cite{jiang2024spatio} approximates arbitrary nonlinearities but requires hundreds of timesteps; and SpikedAttention~\cite{hwang2024spikedattention} achieves fully spike-driven conversion via trace-driven matrix multiplication. 
Despite these advances, all existing methods rely on iterative accumulation and gradual dissipation of residual membrane potential, encoding redundant information across timesteps — a fundamental bottleneck that neither algorithmic nor hardware optimizations have yet resolved.

In contrast, PRIME belongs to the training-free category and requires no fine-tuning or retraining: the conversion is performed in a single forward pass on the pretrained weights, directly parsing the binary representation of floating-point data to generate spikes.